%% Created by Maple 2022.1, Windows 10
%% Source Worksheet: lab_sols11
%% Generated: Wed May 08 10:42:24 BST 2024
\documentclass{article}
\usepackage{amssymb}
\usepackage{graphicx}
\usepackage{hyperref}
\usepackage{listings}
\usepackage{mathtools}
\usepackage{maple}
\usepackage[utf8]{inputenc}
\usepackage[svgnames]{xcolor}
\usepackage{amsmath}
\usepackage{breqn}
\usepackage{textcomp}
\begin{document}
\lstset{basicstyle=\ttfamily,breaklines=true,columns=flexible}
\pagestyle{empty}
\DefineParaStyle{Maple Bullet Item}
\DefineParaStyle{Maple Heading 1}
\DefineParaStyle{Maple Warning}
\DefineParaStyle{Maple Heading 4}
\DefineParaStyle{Maple Heading 2}
\DefineParaStyle{Maple Heading 3}
\DefineParaStyle{Maple Dash Item}
\DefineParaStyle{Maple Error}
\DefineParaStyle{Maple Title}
\DefineParaStyle{Maple Ordered List 1}
\DefineParaStyle{Maple Text Output}
\DefineParaStyle{Maple Ordered List 2}
\DefineParaStyle{Maple Ordered List 3}
\DefineParaStyle{Maple Normal}
\DefineParaStyle{Maple Ordered List 4}
\DefineParaStyle{Maple Ordered List 5}
\DefineCharStyle{Maple 2D Output}
\DefineCharStyle{Maple 2D Input}
\DefineCharStyle{Maple Maple Input}
\DefineCharStyle{Maple 2D Math}
\DefineCharStyle{Maple Hyperlink}
\begin{center}
\begin{Maple Normal}
\underline{\textbf{Miscellaneous problems}}
\end{Maple Normal}
\end{center}
\begin{Maple Normal}
Q1.1
\end{Maple Normal}
\begin{lstlisting}
> with(plots):
\end{lstlisting}
\begin{lstlisting}
> display(
 plot([cos(t)/2-cos(2*t)/4,sin(t)/2-sin(2*t)/4,t=0..2*Pi]),
 plot([cos(t)/4-1,sin(t)/4,t=0..2*Pi]),
 plot([-0.1225+0.0945*cos(t), 0.7449+0.0945*sin(t),t=0..2*Pi]),
 plot([-0.1225+0.0945*cos(t),-0.7449+0.0945*sin(t),t=0..2*Pi]),
 scaling=constrained,axes=none
);
\end{lstlisting}
\begin{Maple Normal}
Here is a picture of the full Mandelbrot set:\

   \

                                   \raisebox{-0.1in}{\includegraphics[keepaspectratio,scale=0.37833333333333335,width=3.1527777777777777in]{lab_sols11_image1.png}}
\end{Maple Normal}
\begin{lstlisting}
> 
\end{lstlisting}
\begin{Maple Normal}

\end{Maple Normal}
\begin{Maple Normal}
\_\_\_\_\_\_\_\_\_\_\_\_\_\_\_\_\_\_\_\_\_\_\_\_\_\_\_\_\_\_\_\_\_\_\_\_\_\_\_\_\_\_\_\_\_\_\_\_\_\_\_\_\_\_\_\_\_\_\_\_\_\_\_\_\_\_\_\_\_\_\_\_\_\_\_\_\_\_\_\_\_
\end{Maple Normal}
\begin{Maple Normal}

\end{Maple Normal}
\begin{Maple Normal}
Q1.2:
\end{Maple Normal}
\begin{lstlisting}
> with(plots):
\end{lstlisting}
\begin{Maple Normal}
Here is the picture for 
$n =6$:
\end{Maple Normal}
\begin{lstlisting}
> implicitplot(abs(x)^6+abs(y)^6=6^(6/2),x=-3..3,y=-3..3,grid=[200,200]);
\end{lstlisting}
\begin{Maple Normal}
Here are 
$n =2$ and 
$n =6$ together:
\end{Maple Normal}
\begin{lstlisting}
> display(
 implicitplot(abs(x)^2+abs(y)^2=2^(2/2),x=-3..3,y=-3..3,grid=[200,200]),
 implicitplot(abs(x)^6+abs(y)^6=6^(6/2),x=-3..3,y=-3..3,grid=[200,200])
);
\end{lstlisting}
\begin{Maple Normal}
Here is the full picture:
\end{Maple Normal}
\begin{lstlisting}
> display(
 seq(
  implicitplot(abs(x)^n+abs(y)^n=n^(n/2),
               x=-3..3,y=-3..3,grid=[200,200]),
  n=1..9)
);
\end{lstlisting}
\begin{lstlisting}
> 

\end{lstlisting}
\begin{Maple Normal}
\_\_\_\_\_\_\_\_\_\_\_\_\_\_\_\_\_\_\_\_\_\_\_\_\_\_\_\_\_\_\_\_\_\_\_\_\_\_\_\_\_\_\_\_\_\_\_\_\_\_\_\_\_\_\_\_\_\_\_\_\_\_\_\_\_\_\_\_\_\_\_\_\_\_\_\_\_\_\_\_\_
\end{Maple Normal}
\begin{Maple Normal}

\end{Maple Normal}
\begin{Maple Normal}
Q1.3:
\end{Maple Normal}
\begin{lstlisting}
> 
\end{lstlisting}
\begin{lstlisting}
> listplot([seq(50^n/n!,n=0..100)],style=point);
\end{lstlisting}
\begin{Maple Normal}
\_\_\_\_\_\_\_\_\_\_\_\_\_\_\_\_\_\_\_\_\_\_\_\_\_\_\_\_\_\_\_\_\_\_\_\_\_\_\_\_\_\_\_\_\_\_\_\_\_\_\_\_\_\_\_\_\_\_\_\_\_\_\_\_\_\_\_\_\_\_\_\_\_\_\_\_\_\_\_\_\_
\end{Maple Normal}
\begin{Maple Normal}
Q1.4
\end{Maple Normal}
\begin{lstlisting}
> with(MTM):
\end{lstlisting}
\begin{lstlisting}
> c := 0.5;
f := (t) -> Re(zeta(c+I*t)):
g := (t) -> Im(zeta(c+I*t)):
\end{lstlisting}
\begin{lstlisting}
> plot([f(t),g(t)],t=0..15);
\end{lstlisting}
\begin{lstlisting}
> plot([f(t),g(t)],t=14..14.3);
\end{lstlisting}
\begin{lstlisting}
> plot([f(t),g(t),t=0..50]);
\end{lstlisting}
\begin{lstlisting}
> 
\end{lstlisting}
\begin{lstlisting}
> 
\end{lstlisting}
\begin{lstlisting}
> 
\end{lstlisting}
\begin{Maple Normal}
\_\_\_\_\_\_\_\_\_\_\_\_\_\_\_\_\_\_\_\_\_\_\_\_\_\_\_\_\_\_\_\_\_\_\_\_\_\_\_\_\_\_\_\_\_\_\_\_\_\_\_\_\_\_\_\_\_\_\_\_\_\_\_\_\_\_\_\_\_\_\_\_\_\_\_\_\_\_\_\_\_
\end{Maple Normal}
\begin{Maple Normal}
Q1.4:
\end{Maple Normal}
\begin{lstlisting}
> y := exp(-1/x);
\end{lstlisting}
\begin{Maple Normal}
We can work out the first few derivatives as follows:
\end{Maple Normal}
\begin{lstlisting}
> diff(y,x);
\end{lstlisting}
\begin{lstlisting}
> simplify(diff(y,x$2));
\end{lstlisting}
\begin{lstlisting}
> simplify(diff(y,x$3));
\end{lstlisting}
\begin{lstlisting}
> simplify(diff(y,x$4));
\end{lstlisting}
\begin{lstlisting}
> simplify(diff(y,x$5));
\end{lstlisting}
\begin{Maple Normal}
We see that 
$\frac{{\textcolor{gray}{\partial}}^{k}}{\textcolor{gray}{\partial}x^{k}} y$ is always of the form 
$\frac{{\mathrm e}^{-\frac{1}{x}} p_{k}(x)}{x^{2 k}}$ for some polynomial 
$p_{k}(x)$.  Explicitly, the first ten of these \

polynomials are as follows:\


\end{Maple Normal}
\begin{lstlisting}
> for k from 1 to 10 do p[k] := expand(x^(2*k)*diff(y,x$k)/y); od;
\end{lstlisting}
\begin{Maple Normal}
We see from this that 
$p_{k}(x)$ always has degree 
$k -1$.  The constant term is always 
$p_{k}(0)=1$.  You should recognise \

the numbers 
$2,6,24,120,720$ as factorials, so the highest term in 
$p_{k}(x)$ is 
$(-1)^{k} k ! x^{k -1}$.  
\end{Maple Normal}
\begin{lstlisting}
> 
\end{lstlisting}
\begin{Maple Normal}
We can plot the derivatives up to 
$k =4$ as follows.  We have divided each function by its approximate\

maximum value so that we can see them clearly in the same picture.   We find that the higher derivatives\

oscillate wildly for moderately small values of 
$x$, but then flatten out for very small values of 
$x$, and also\

for reasonably large values of 
$x$..
\end{Maple Normal}
\begin{lstlisting}
> plot([
y/0.6,
diff(y,x)/0.55,
diff(y,x,x)/2.5,
diff(y,x,x,x)/30,
diff(y,x,x,x,x)/600],
x=0..2);
\end{lstlisting}
\begin{Maple Normal}
This plot of the sixth derivative reinforces the same message.
\end{Maple Normal}
\begin{lstlisting}
> plot(diff(y,x$6),x=0..3,-100..100);
\end{lstlisting}
\begin{Maple Normal}
We can watch the curves flattening out near the origin as follows:
\end{Maple Normal}
\begin{lstlisting}
> plot([
y,
diff(y,x),
diff(y,x,x),
diff(y,x,x,x),
diff(y,x,x,x,x)],
x=0..0.1,0..0.002);
\end{lstlisting}
\begin{lstlisting}
> 
\end{lstlisting}
\begin{Maple Normal}
\_\_\_\_\_\_\_\_\_\_\_\_\_\_\_\_\_\_\_\_\_\_\_\_\_\_\_\_\_\_\_\_\_\_\_\_\_\_\_\_\_\_\_\_\_\_\_\_\_\_\_\_\_\_\_\_\_\_\_\_\_\_\_\_\_\_\_\_\_\_\_\_\_\_\_\_\_\_\_\_\_
\end{Maple Normal}
\begin{Maple Normal}

\end{Maple Normal}
\begin{center}
\begin{Maple Normal}
\underline{\textbf{Solitons}}
\end{Maple Normal}
\end{center}
\begin{Maple Normal}
\_\_\_\_\_\_\_\_\_\_\_\_\_\_\_\_\_\_\_\_\_\_\_\_\_\_\_\_\_\_\_\_\_\_\_\_\_\_\_\_\_\_\_\_\_\_\_\_\_\_\_\_\_\_\_\_\_\_\_\_\_\_\_\_\_\_\_\_\_\_\_\_\_\_\_\_\_\_\_\_\_
\end{Maple Normal}
\begin{lstlisting}
> with(plots):
\end{lstlisting}
\begin{Maple Normal}
We can enter the basic definitions as follows.\


\end{Maple Normal}
\begin{lstlisting}
> q := sqrt(2);
p := log(3 + 2*q);
r := x-4*t;
s := q*(x - 8*t);
T := 32 * cosh(2*r-p) + 16 * cosh(2*s-p) + 16;
B := 4*(1+q) * cosh(r)*cosh(s) + (4*q-8)*exp(r+s);
phi[0] := 2*cosh(r)^(-2);
phi[1] := 4*cosh(s)^(-2);
phi[2] := 2*cosh(r-p)^(-2);
phi[3] := 4*cosh(s-p)^(-2);
phi[4] := T/B^2;
\end{lstlisting}
\begin{lstlisting}
> 
\end{lstlisting}
\begin{Maple Normal}

\end{Maple Normal}
\begin{Maple Normal}
We now check that the Korteweg-de Vries equation is satisfied.  The tidiest way is to introduce the KdV operator as follows:\


\end{Maple Normal}
\begin{lstlisting}
> K := (u) -> diff(u,t) + diff(u,x,x,x) + 6 * u * diff(u,x);
\end{lstlisting}
\begin{Maple Normal}
We now apply 
$K$ to the functions 
$\varphi_{i}$:
\end{Maple Normal}
\begin{lstlisting}
> simplify(K(phi[0]));
simplify(K(phi[1]));
simplify(K(phi[2]));
simplify(K(phi[3]));
\end{lstlisting}
\begin{Maple Normal}
If we just apply 
$K$ directly to 
$\varphi_{4}$ then we get several pages of dense output, because Maple is not very good at simplifying\

hyperbolic functions.  As explained in the problem sheet, we need to convert everything to exponential form first.
\end{Maple Normal}
\begin{lstlisting}
> convert(phi[4],exp);
\end{lstlisting}
\begin{lstlisting}
> 
\end{lstlisting}
\begin{lstlisting}
> simplify(K(convert(phi[4],exp)));
\end{lstlisting}
\begin{lstlisting}
> 
\end{lstlisting}
\begin{Maple Normal}
We can plot all the functions 
$\varphi_{i}$ together as follows.  The bottom one is 
$\varphi_{0}$ and the top one is 
$\varphi_{4}$.
\end{Maple Normal}
\begin{lstlisting}
> animate(
plot,
[[phi[0],phi[1]+5,phi[2]+10,phi[3]+15,phi[4]+20],x=-30..30],
t=-3..3,
frames=50,
scaling=constrained,
axes=none);
\end{lstlisting}
\begin{Maple Normal}
We see that 
$\varphi_{0}$ is essentially the same as 
$\varphi_{2}$ but delayed slightly; this is already easy to see from the formulae.\

Similarly, 
$\varphi_{1}$ is a slightly delayed copy of 
$\varphi_{3}$.  When 
$t$ is negative, the small hump in 
$\varphi_{4}$ follows 
$\varphi_{2}$ and the \

large hump in 
$\varphi_{4}$ follows 
$\varphi_{1}$, so 
$\varphi_{4}$ is approximately 
$\varphi_{1}+\varphi_{2}$.  Near 
$t =0$ the two humps interact, the large one\

jumps forward a little to follow 
$\varphi_{3}$, and the small hump drops back a little to follow 
$\varphi_{0}$.  Thus, when 
$t$ is positive\

we see that 
$\varphi_{4}$ is approximately 
$\varphi_{0}+\varphi_{3}$.\

\

We can see all this again in the following animation.  The middle graph is 
$\varphi_{4}$.  The bottom graph is 
$\varphi_{4}-\varphi_{1}-\varphi_{2}$; \

we find that this is very small when 
$t$ is negative.  The top graph is 
$\varphi_{4}-\varphi_{0}-\varphi_{3}$; we find that this is very small when\


$t$ is positive. 
\end{Maple Normal}
\begin{lstlisting}
> animate(
plot,
[[phi[4]-phi[1]-phi[2]-10,phi[4],phi[4]-phi[0]-phi[3]+10],x=-30..30],
t=-3..3,
frames=50,
scaling=constrained,
axes=none);
\end{lstlisting}
\begin{Maple Normal}
We now calculate the momenta of 
$\varphi_{0}..\varphi_{3}$:
\end{Maple Normal}
\begin{lstlisting}
> M[0] := int(phi[0],x=-infinity..infinity);
M[1] := int(phi[1],x=-infinity..infinity);
M[2] := int(phi[2],x=-infinity..infinity);
M[3] := int(phi[3],x=-infinity..infinity);
\end{lstlisting}
\begin{Maple Normal}
We can do this more efficiently as follows:
\end{Maple Normal}
\begin{lstlisting}
> for i from 0 to 3 do
 M[i] := int(phi[i],x=-infinity..infinity);
od;
\end{lstlisting}
\begin{Maple Normal}
For 
$M_{4}$ we need a more elaborate method, as described in the problem sheet:
\end{Maple Normal}
\begin{lstlisting}
> M[4] := int(subs(t=0,phi[4]),x=-infinity..infinity,numeric,method=_Gquad);
\end{lstlisting}
\begin{Maple Normal}
We find that 
$M_{4}=M_{1}+M_{2}$, apart from a tiny error that would go away if we computed the integral more accurately.
\end{Maple Normal}
\begin{lstlisting}
> evalf(M[4] - M[1] - M[2]);
\end{lstlisting}
\begin{lstlisting}
> 
\end{lstlisting}
\begin{Maple Normal}
We now repeat this for the energy:
\end{Maple Normal}
\begin{lstlisting}
> for i from 0 to 3 do
 E[i] := int(phi[i]^2,x=-infinity..infinity);
od;
\end{lstlisting}
\begin{lstlisting}
> 
\end{lstlisting}
\begin{lstlisting}
> E[4] := int(subs(t=0,phi[4]^2),x=-infinity..infinity,numeric,method=_Gquad);
\end{lstlisting}
\begin{lstlisting}
> evalf(E[4] - E[1] - E[2]);
\end{lstlisting}
\begin{lstlisting}
> 
\end{lstlisting}
\begin{lstlisting}
> 
\end{lstlisting}
\end{document}
